\documentclass[12pt]{article}
\usepackage[utf8]{inputenc}
\usepackage[spanish]{babel}
\usepackage{amsmath}
\usepackage{amsfonts}
\usepackage{amssymb}
\usepackage{geometry}
\geometry{letterpaper, margin=2.5cm}

\begin{document}

\begin{center}
    \Large \textbf{Evaluación de Examen 1: Unidad I} \\
    \large Física para Ciencias de la Salud (005-1713) \\
    \vspace{0.5cm}
    \textbf{Estudiante:} Krischel Bravo \quad \textbf{C.I.:} 32.449.538
    \vspace{0.3cm} \\
    \large \textbf{Calificación Final: 9/10 puntos}
\end{center}

\vspace{0.5cm}

\section*{Análisis de Resultados}

\subsection*{1. Ángulo plano (Conversión a radianes)}
\textbf{Procedimiento y resultado:} Plantea correctamente la conversión usando el factor equivalente a $180^\circ$. A pesar de un leve error de transcripción inicial (escribió 18 en vez de 180), lo corrige en el acto y realiza una simplificación de fracciones muy limpia $\left(\frac{75\pi}{180} = \frac{15}{36} = \frac{5\pi}{12}\right)$, obteniendo el valor exacto de $\frac{5\pi}{12}$ rad. \\
\textbf{Veredicto:} \textbf{Correcto} (2/2 pts).

\subsection*{2. Sistemas de coordenadas (Longitud del hueso)}
\textbf{Procedimiento y resultado:} Utiliza la ecuación formal de la distancia entre dos puntos. Sustituye correctamente las diferencias en cada eje y eleva al cuadrado $\sqrt{6^2 + 8^2}$. Finalmente aplica la raíz cuadrada ($\sqrt{100}$) y llega exitosamente a $10$ cm. \\
\textbf{Veredicto:} \textbf{Correcto} (2/2 pts).

\subsection*{3. Vectores (Fuerza resultante)}
\textbf{Procedimiento y resultado:} La estudiante es muy ordenada al desglosar el cálculo por componentes: resuelve primero la sumatoria en el eje $x$ ($45\hat{i} + (-15\hat{i}) = 30\hat{i}$) y luego en el eje $y$ ($25\hat{j} + 35\hat{j} = 60\hat{j}$). Ensambla exitosamente el vector resultante $\vec{F}_R = 30\hat{i} + 60\hat{j}$, agregando de manera un poco informal la unidad final (= N). \\
\textbf{Veredicto:} \textbf{Correcto} (2/2 pts).

\subsection*{4. Potencias de diez (Notación científica y prefijo)}
\textbf{Procedimiento y resultado:} Expresa la medida decimal de forma adecuada empleando una potencia de base diez ($12,5 \cdot 10^{-6}$ m). Sin embargo, omite la respuesta final específica requerida en el enunciado, la cual consistía en identificar e indicar explícitamente el nombre del prefijo del Sistema Internacional (micrómetros). \\
\textbf{Veredicto:} \textbf{Incompleto / Parcialmente correcto} (1/2 pts).

\subsection*{5. Sistemas de unidades (Conversión de densidad)}
\textbf{Procedimiento y resultado:} La estudiante logra deducir que el factor global de conversión entre $\text{g/cm}^3$ y $\text{kg/m}^3$ equivale a multiplicar por $1000$. Ejecuta la operación aritmética ($1,03 \times 1000$) y obtiene el resultado final correcto ($1030$ de forma exacta). No obstante, el procedimiento analítico presentado en la hoja es bastante desordenado y confuso. \\
\textbf{Veredicto:} \textbf{Correcto} (2/2 pts).

\vspace{0.5cm}
\section*{Comentarios Generales y Sugerencias}
El examen es muy bueno. La estudiante demuestra que comprende a la perfección la lógica detrás de los ejercicios de matemáticas y álgebra vectorial.
\begin{itemize}
    \item Se recomienda enfáticamente cuidar el orden y la formalidad al escribir los procedimientos matemáticos (como en la Pregunta 5), ya que estructurar mejor los pasos con factores de conversión en cadena evitará confusiones en el futuro.
    \item Como sucedió con varios compañeros, faltó leer detalladamente la instrucción en la Pregunta 4 para incluir el nombre del prefijo.
\end{itemize}

\end{document}